\section{锅炉}

%! src=8 kind=orig title=锅炉·章节导入
\begin{frame}[t,shrink]{\src{8}{锅炉：章节导入}}
\chapterintro{02}{锅炉}{Steam generation: combustion, heat transfer and heat balance}
{\begin{itemize}
  \item 锅炉产生什么参数、什么流量的蒸汽？
  \item “炉”和“锅”分别承担什么过程？
  \item 燃料热量怎样分配？效率如何定义和测量？
\end{itemize}}
{原页 8--24 \,$\cdot$\, 主线：燃料燃烧 $\rightarrow$ 烟气放热 $\rightarrow$ 工质吸热 $\rightarrow$ 蒸汽}
\vspace{3mm}
\begin{center}
\small
\textbf{关键词}\quad 燃烧 \enspace$\cdot$\enspace 受热面 \enspace$\cdot$\enspace 水循环 \enspace$\cdot$\enspace 热平衡 \enspace$\cdot$\enspace 煤质基准
\end{center}
\end{frame}

%! src=9 kind=orig title=锅炉的作用与蒸汽参数
\begin{frame}[t,shrink]{\src{9}{锅炉的作用与蒸汽参数}}
\begin{columns}[T,onlytextwidth]
\begin{column}{0.56\textwidth}
  \begin{itemize}
    \item 锅炉完成循环工质的\textbf{加热、汽化、蒸汽过热}。
    \item 能量转换：燃料化学能 $\rightarrow$ 烟气热能 $\rightarrow$ 工质焓增加。
    \item 蒸汽的“质”：压力 \(p\)、温度 \(t\)、纯净度。
    \item 蒸汽的“量”：蒸发量 \(D\)，单位时间产生的蒸汽质量流量。
    \item 三个耦合过程：\textbf{燃料燃烧、工质吸热、烟气放热}。
    \item 锅炉还可承担汽轮机高压缸排汽的再热。
  \end{itemize}
\end{column}
\begin{column}{0.42\textwidth}
  \begin{keybox}{“质”和“量”要同时给出}
  \small 只写“产生蒸汽”不足以描述锅炉。工程上至少给出压力、温度、流量和给水条件；再热机组还要给出再热蒸汽的参数与流量。
  \end{keybox}
  \vspace{0.8mm}
  \begin{plainbox}
  \small
  \(D=\dot m_{\mathrm{s}}\)：主蒸汽质量流量。\\
  常用工程单位：\(\mathrm{t/h}\)；SI 单位：\(\mathrm{kg/s}\)。
  \end{plainbox}
  \vspace{0.8mm}
  \begin{amberbox}{不要混淆}
  \small 蒸发量是\textbf{流量}，不是热量，也不是锅炉容量本身；同样的 \(D\)，参数越高，所需吸热量越大。
  \end{amberbox}
\end{column}
\end{columns}
\end{frame}

%! src=9 kind=supp title=蒸发量怎样联系到有效吸热
\begin{frame}[t,shrink]{\srcsupp{9}{蒸发量怎样联系到锅炉有效吸热}}
稳态、忽略排污等支流时，主蒸汽与再热蒸汽的有效吸热率可写为：
$$
\dot Q_1
=\dot m_{\mathrm{ms}}\left(h_{\mathrm{ms}}-h_{\mathrm{fw}}\right)
+\dot m_{\mathrm{rh}}\left(h_{\mathrm{rh,out}}-h_{\mathrm{rh,in}}\right)
$$
\begin{columns}[T,onlytextwidth]
\begin{column}{0.52\textwidth}
  \begin{itemize}
    \item \(\dot m_{\mathrm{ms}}\)、\(\dot m_{\mathrm{rh}}\)：主蒸汽、再热蒸汽质量流量，\(\mathrm{kg/s}\)。
    \item \(h_{\mathrm{ms}}\)、\(h_{\mathrm{fw}}\)：主蒸汽、给水比焓，\(\mathrm{kJ/kg}\)。
    \item \(h_{\mathrm{rh,out}}\)、\(h_{\mathrm{rh,in}}\)：再热器出口、入口比焓。
    \item \(\dot Q_1\) 的单位为 \(\mathrm{kW}\)；若用 \(\mathrm{kJ/h}\)，则需另换单位。
  \end{itemize}
\end{column}
\begin{column}{0.45\textwidth}
  \begin{keybox}{单位换算}
  \small
  \(1\,\mathrm{t/h}=1000/3600\,\mathrm{kg/s}=1/3.6\,\mathrm{kg/s}\)。
  \par\vspace{1mm}
  \(1\,\mathrm{kW}=1\,\mathrm{kJ/s}=3600\,\mathrm{kJ/h}\)。
  \end{keybox}
  \vspace{0.7mm}
  \begin{amberbox}{常见错误}
  \small 用 \(D\) 直接代替 \(\dot Q_1\)，或用 \(c_p\Delta T\) 跨过汽化区计算过热蒸汽吸热量。
  \end{amberbox}
\end{column}
\end{columns}
\end{frame}

%! src=10 kind=orig title=锅炉的两类分类依据
\begin{frame}[t,shrink]{\src{10}{锅炉的两类分类依据}}
\begin{columns}[T,onlytextwidth]
\begin{column}{0.49\textwidth}
  \begin{keybox}{按水循环方式}
  \small
  \begin{tabularx}{\linewidth}{@{}>{\bfseries}p{14mm}X@{}}
    自然循环 & 利用下降管与上升管的密度差\\
    强制循环 & 增加循环泵提供压头\\
    直流 & 给水泵压头使工质一次通过\\
  \end{tabularx}
  \end{keybox}
  \vspace{1mm}
  \begin{keybox}{按出口蒸汽参数}
  \small 低压、中压、高压、超高压、亚临界、超临界、超超临界。
  \end{keybox}
\end{column}
\begin{column}{0.48\textwidth}
  \begin{itemize}
    \item \textbf{水循环方式}描述水/汽在受热面中如何流动，以及有无汽包。
    \item \textbf{蒸汽参数}描述出口蒸汽所处的压力、温度等级。
    \item 两者是不同维度：超临界锅炉通常采用直流，但“超临界”不等于“直流”的定义。
    \item 临界参数（课件）：\(p_{\mathrm{cr}}=22.129\,\mathrm{MPa}\)，\(t_{\mathrm{cr}}=374.15\,\mathrm{^\circ C}\)。
  \end{itemize}
  \vspace{0.5mm}
  \begin{amberbox}{易混点}
  \small “超超临界”是工程参数等级，不是热力学上另一个临界点；它仍以水的临界参数为参照。
  \end{amberbox}
\end{column}
\end{columns}
\end{frame}

%! src=11 kind=orig title=“炉”与“锅”的结构分工
\begin{frame}[t,shrink]{\src{11}{“炉”与“锅”的结构分工}}
\begin{columns}[T,onlytextwidth]
\begin{column}{0.48\textwidth}
  \begin{itemize}
    \item \textbf{炉侧}：燃料在炉膛中燃烧，化学能变为高温烟气热能。
    \item \textbf{锅侧}：工质经过预热、蒸发、过热三个阶段吸热。
    \item 省煤器：未饱和水 $\rightarrow$ 接近或达到饱和的水。
    \item 水冷壁：水受热形成汽水混合物。
    \item 汽包：汽水分离、汇集并组织循环。
    \item 过热器：饱和蒸汽 $\rightarrow$ 过热蒸汽。
  \end{itemize}
\end{column}
\begin{column}{0.50\textwidth}
  \begin{amberbox}{不要把汽包当成“蒸发受热面”}
  \small 汽包主要承担汽水分离、储水和连接循环回路；真正的蒸发热量主要在水冷壁等蒸发受热面中吸收。
  \end{amberbox}
  \vspace{0.8mm}
  \begin{keybox}{读结构图}
  \small 先分“烟气侧”和“工质侧”，再看两者在哪些受热面交换热量。烟气与水/蒸汽通常不直接混合。
  \end{keybox}
\end{column}
\end{columns}
\vspace{1mm}
\begin{keybox}{结构分工一句话}
\small 炉侧：燃料 + 空气 $\rightarrow$ 烟气；锅侧：给水 $\rightarrow$ 蒸汽。两侧只通过受热面交换热量，烟气与水/蒸汽不直接混合。
\end{keybox}
\end{frame}

%! src=11 kind=supp title=预热、蒸发、过热的热量
\begin{frame}[t,shrink]{\srcsupp{11}{预热、蒸发、过热的热量分别从哪里来}}
同一亚临界压力下，忽略压降，以单位质量工质为基准：
$$
\begin{aligned}
q_{\mathrm{pre}} &= h_{\mathrm{f}}-h_{\mathrm{fw}},\\
q_{\mathrm{evap}} &= h_{\mathrm{g}}-h_{\mathrm{f}},\\
q_{\mathrm{sh}} &= h_{\mathrm{s}}-h_{\mathrm{g}}
\end{aligned}
$$
\begin{columns}[T,onlytextwidth]
\begin{column}{0.49\textwidth}
  \begin{itemize}
    \item \(h_{\mathrm{f}}\)、\(h_{\mathrm{g}}\)：该压力下饱和水、饱和蒸汽比焓。
    \item 预热主要升温；蒸发主要改变相态；过热继续提高温度。
    \item 三段热量由不同受热面承担，但烟气侧并不按“三段”截然分开。
  \end{itemize}
\end{column}
\begin{column}{0.48\textwidth}
  \begin{amberbox}{超临界时怎么理解}
  \small 超临界压力下不存在通常意义的两相沸腾平台，不宜硬套 \(h_{\mathrm{g}}-h_{\mathrm{f}}\)。应直接用压力、温度查取物性，按焓差计算。
  \end{amberbox}
  \vspace{0.8mm}
  \begin{keybox}{公式条件}
  \small 稳态流动、忽略动能和位能变化；若进出口速度很大，应改用总焓。
  \end{keybox}
\end{column}
\end{columns}
\end{frame}

%! src=12 kind=orig title=锅炉本体的主要部件
\begin{frame}[t,shrink]{\src{12}{锅炉本体的主要部件}}
\begin{center}
\begin{tikzpicture}[x=1mm,y=1mm,scale=0.9,transform shape]
  \node[box,text width=28mm,minimum height=10mm] (a) at (18,0) {炉膛 / 燃烧器\\燃烧放热};
  \node[box,text width=43mm,minimum height=10mm] (b) at (70,0) {受热面组\\水冷壁 · 过热器 · 再热器 · 省煤器 · 空预器};
  \node[boxb,text width=28mm,minimum height=10mm] (c) at (124,0) {工质吸热\\给水 $\rightarrow$ 蒸汽};
  \draw[arr] (a) -- node[lab,above]{高温烟气} (b);
  \draw[arr] (b) -- node[lab,above]{热量传递} (c);
  \node[lab,text=cprimary!80!black] at (70,-10) {燃烧空间 $\rightarrow$ 受热面 $\rightarrow$ 工质};
\end{tikzpicture}
\end{center}
\begin{smalltab}
\begin{tabularx}{\linewidth}{@{}>{\bfseries}p{20mm}X>{\bfseries}p{20mm}X@{}}
  \toprule
  部件 & 主要作用 & 部件 & 主要作用\\
  \midrule
  炉膛、燃烧器 & 组织燃烧，形成高温烟气 & 水冷壁 & 吸收炉膛辐射热，使水蒸发\\
  汽包 & 汽水分离、储水、连接循环 & 过热器 & 把饱和蒸汽加热为过热蒸汽\\
  再热器 & 加热汽轮机高压缸排汽 & 省煤器 & 用烟气余热预热给水\\
  空气预热器 & 用烟气余热加热助燃空气 & 炉墙 & 保温、减少散热\\
  \bottomrule
\end{tabularx}
\end{smalltab}
\vspace{0.5mm}
\begin{keybox}{记忆线索}
\small 燃烧空间 + 汽水受热面 + 尾部余热利用设备。
\end{keybox}
\end{frame}

%! src=13 kind=orig title=锅炉辅助设备
\begin{frame}[t,shrink]{\src{13}{锅炉辅助设备}}
课件列出八类辅助设备或功能。按“介质链”分组更容易记：
\begin{columns}[T,onlytextwidth]
\begin{column}{0.49\textwidth}
  \begin{keybox}{燃料与风烟}
  \small
  \textbf{燃料供应}、\textbf{磨煤机}：供煤并制备煤粉。\\
  \textbf{送风}：提供助燃空气。\\
  \textbf{引风}：克服烟道阻力，维持炉膛负压。\\
  \textbf{除灰除渣}：排出飞灰和底渣。
  \end{keybox}
  \vspace{0.8mm}
  \begin{keybox}{汽水与环保}
  \small
  \textbf{给水}：向锅炉输送合格给水。\\
  \textbf{脱硫脱硝}：控制 \(\mathrm{SO_2}\)、\(\mathrm{NO_x}\)。\\
  \textbf{自动控制}：协调燃烧、汽水参数与辅机。
  \end{keybox}
\end{column}
\begin{column}{0.48\textwidth}
  \begin{itemize}
    \item 本体直接完成燃烧与传热；辅助设备保证本体连续、稳定、清洁运行。
    \item 送风机与引风机作用不同：一个送空气，一个抽烟气。
    \item 磨煤机同时影响煤粉细度、干燥和输送；不是简单的“粉碎”。
    \item 除灰除渣与脱硫脱硝处理的是不同形态的污染物。
  \end{itemize}
  \vspace{0.5mm}
  \begin{amberbox}{运行联系}
  \small 风量、煤量、给水量和引风量必须协调；任一环节波动都会反映到蒸汽参数、炉膛负压和排烟损失上。
  \end{amberbox}
\end{column}
\end{columns}
\end{frame}

%! src=14 kind=orig title=按蒸汽参数比较机组
\begin{frame}[t,shrink]{\src{14}{按蒸汽参数比较机组}}
课件表摘录（压力和温度单位分别为 \(\mathrm{MPa}\)、\(\mathrm{^\circ C}\)；供电煤耗单位 \(\mathrm{g/(kW\cdot h)}\)）：
\begin{center}
\begin{tabularx}{0.92\linewidth}{@{}lcccc@{}}
  \toprule
  类型 & 压力 & 温度 & 电厂效率 & 供电煤耗\\
  \midrule
  亚临界 & 17.0 & 540/540 & \(38\%\) & 324\\
  超临界 & 25.5 & 567/567 & \(41\%\) & 300\\
  超超临界 & 30.0 & 600/600/600 & \(48\%\) & 256\\
  \bottomrule
\end{tabularx}
\end{center}
\begin{columns}[T,onlytextwidth]
\begin{column}{0.52\textwidth}
  \begin{itemize}
    \item 课件趋势：提高蒸汽参数，有利于提高效率、降低煤耗。
    \item 临界参数：\(p_{\mathrm{cr}}=22.129\,\mathrm{MPa}\)，\(t_{\mathrm{cr}}=374.15\,\mathrm{^\circ C}\)。
    \item 表中效率、煤耗是课件对比示例，不能直接外推为所有机组。
  \end{itemize}
\end{column}
\begin{column}{0.45\textwidth}
  \begin{amberbox}{读表注意}
  \small 温度栏用斜线分隔多个值，但课件未逐项注明是主蒸汽、一次再热还是二次再热。考试按课件表记忆；工程比较时应查清参数定义。
  \end{amberbox}
\end{column}
\end{columns}
\end{frame}

%! src=15 kind=orig title=自然循环、强制循环与直流锅炉
\begin{frame}[t,shrink]{\src{15}{自然循环、强制循环与直流锅炉}}
\begin{columns}[T,onlytextwidth]
\begin{column}{0.32\textwidth}
  \begin{tikzpicture}[x=1mm,y=1mm,scale=0.82,transform shape]
    \draw[cprimary,thick] (3,4) rectangle (33,30);
    \draw[cprimary,thick] (16,30) circle (4);
    \draw[cprimary,thick] (8,30) -- (8,8) -- (28,8) -- (28,30);
    \draw[arr] (8,18) -- (8,8);
    \draw[arr] (28,10) -- (28,28);
    \node[lab,anchor=south] at (16,35) {汽包};
    \node[lab,rotate=90] at (5,19) {下降管};
    \node[lab,rotate=90] at (31,19) {上升管};
    \node[lab] at (18,4) {自然循环};
  \end{tikzpicture}
\end{column}
\begin{column}{0.32\textwidth}
  \begin{tikzpicture}[x=1mm,y=1mm,scale=0.82,transform shape]
    \draw[cprimary,thick] (3,4) rectangle (33,30);
    \draw[cprimary,thick] (16,30) circle (4);
    \draw[cprimary,thick] (8,30) -- (8,8) -- (28,8) -- (28,30);
    \draw[arr] (8,18) -- (8,8);
    \draw[arr] (28,10) -- (28,28);
    \draw[fill=white,thick] (8,4) circle (3);
    \node[lab] at (8,0) {循环泵};
    \node[lab] at (24,4) {强制循环};
  \end{tikzpicture}
\end{column}
\begin{column}{0.32\textwidth}
  \begin{tikzpicture}[x=1mm,y=1mm,scale=0.82,transform shape]
    \draw[cprimary,thick] (3,4) rectangle (33,30);
    \draw[cprimary,thick] (8,6) -- (8,26) -- (18,26) -- (18,10) -- (28,10) -- (28,30);
    \draw[arr] (8,12) -- (8,24);
    \draw[arr] (18,20) -- (18,12);
    \draw[arr] (28,16) -- (28,28);
    \node[lab] at (18,4) {直流};
    \node[lab,anchor=south] at (28,31) {去过热器};
  \end{tikzpicture}
\end{column}
\end{columns}
\vspace{1mm}
\begin{smalltab}
\begin{tabularx}{\linewidth}{@{}>{\bfseries}p{17mm}p{13mm}p{44mm}X@{}}
  \toprule
  类型 & 汽包 & 主要驱动力 & 流动特点\\
  \midrule
  自然循环 & 有 & 下降管与上升管密度差 & 回路内循环\\
  强制循环 & 有 & 密度差 + 循环泵压头 & 泵辅助循环\\
  直流 & 无 & 给水泵压头 & 一次通过各受热面\\
  \bottomrule
\end{tabularx}
\end{smalltab}
\begin{amberbox}{记忆}
\small 自然循环靠“密度差”，强制循环靠“泵”，直流靠“给水泵一次推过去”。超临界锅炉通常采用直流，因为不存在稳定的汽水两相密度差循环。
\end{amberbox}
\end{frame}

%! src=15 kind=supp title=自然循环的驱动力
\begin{frame}[t,shrink]{\srcsupp{15}{自然循环的驱动力与循环倍率}}
\begin{columns}[T,onlytextwidth]
\begin{column}{0.45\textwidth}
\begin{tikzpicture}[x=1mm,y=1mm,scale=0.82,transform shape]
  \draw[cprimary,thick] (5,5) rectangle (55,45);
  \draw[cprimary,thick] (30,45) circle (5);
  \draw[cprimary,thick] (12,45) -- (12,10) -- (48,10) -- (48,45);
  \draw[arr] (12,30) -- (12,12);
  \draw[arr] (48,12) -- (48,42);
  \node[lab,anchor=south] at (30,51) {汽包};
  \node[lab,rotate=90] at (9,28) {下降管：水};
  \node[lab,rotate=90] at (51,28) {上升管：汽水混合物};
  \draw[<->,thin] (57,10) -- (57,44) node[midway,right,lab] {\(H\)};
\end{tikzpicture}
\end{column}
\begin{column}{0.52\textwidth}
  设下降管密度为 \(\rho_{\mathrm{d}}\)，上升管平均密度为 \(\bar\rho_{\mathrm{r}}\)，忽略沿程加速：
  $$
  \Delta p_{\mathrm{dr}}\approx gH\left(\rho_{\mathrm{d}}-\bar\rho_{\mathrm{r}}\right)
  $$
  \begin{itemize}
    \item 下降管中主要是水，密度较大；上升管中为汽水混合物，平均密度较小。
    \item 该压差必须克服循环回路的流动阻力，才能建立稳定循环。
    \item 循环倍率 \(n=\dot m_{\mathrm{circ}}/\dot m_{\mathrm{steam}}\) 描述循环水量是蒸发量的多少倍。
  \end{itemize}
\end{column}
\end{columns}
\vspace{0.8mm}
\begin{amberbox}{公式条件}
\small 这是忽略加速和局部损失的简化驱动压头；实际循环还受锅炉负荷、含汽率、管路阻力和两相流结构影响。
\end{amberbox}
\end{frame}

%! src=16 kind=orig title=锅炉系统总览
\begin{frame}[t,shrink]{\src{16}{锅炉系统总览}}
课件把锅炉系统分为四部分：
\begin{columns}[T,onlytextwidth]
\begin{column}{0.49\textwidth}
  \begin{keybox}{煤粉制备}
  \small 煤仓、给煤机、磨煤机、送粉设备。任务：供煤、干燥、磨细、输送。
  \end{keybox}
  \vspace{0.7mm}
  \begin{keybox}{燃烧系统}
  \small 燃烧器、炉膛、水冷壁、排渣装置。任务：组织着火、稳定燃烧、冷却炉膛。
  \end{keybox}
\end{column}
\begin{column}{0.49\textwidth}
  \begin{keybox}{风烟系统}
  \small 送风机、空预器、烟道、除尘、引风机、脱硫脱硝、烟囱。任务：供氧、排烟、净化。
  \end{keybox}
  \vspace{0.7mm}
  \begin{keybox}{汽水系统}
  \small 给水、汽包、下降管、水冷壁、过热器、再热器。任务：吸热、蒸发、过热。
  \end{keybox}
\end{column}
\end{columns}
\vspace{1mm}
\begin{plainbox}
\small \textbf{读总图方法}：先按介质分线（煤、空气、烟气、水/蒸汽），再沿流向找设备，最后把每条线与它的能量交换点对应起来；不要从编号开始背。
\end{plainbox}
\end{frame}

%! src=17 kind=orig title=煤粉制备和燃烧流程
\begin{frame}[t,shrink]{\src{17}{煤粉制备和燃烧流程}}
\begin{center}
\begin{tikzpicture}[x=1mm,y=1mm,scale=0.82,transform shape]
  \node[box,text width=20mm] (coal) at (10,10) {煤仓\\原煤};
  \node[box,text width=20mm] (feed) at (38,10) {给煤机};
  \node[box,text width=22mm] (mill) at (67,10) {磨煤机\\煤粉};
  \node[box,text width=20mm] (burn) at (97,10) {燃烧器};
  \node[box,text width=20mm] (furn) at (127,10) {炉膛};
  \node[boxg,text width=20mm] (air) at (67,-8) {热空气};
  \node[boxg,text width=22mm] (ap) at (31,-8) {空气预热器};
  \node[boxg,text width=20mm] (fd) at (4,-8) {送风机};
  \draw[arr] (coal) -- (feed);
  \draw[arr] (feed) -- (mill);
  \draw[arr] (mill) -- (burn);
  \draw[arr] (burn) -- (furn);
  \draw[arr] (fd) -- (ap);
  \draw[arr] (ap) -- (air);
  \draw[arr] (air) -- (mill);
  \draw[arr] (air) -- (burn);
  \node[lab,text=cprimary] at (129,18) {高温烟气};
\end{tikzpicture}
\end{center}
\begin{columns}[T,onlytextwidth]
\begin{column}{0.49\textwidth}
  \begin{itemize}
    \item 磨细煤粒扩大与氧气的接触面积，利于着火和燃尽。
    \item 热空气同时承担\textbf{干燥、输送}和部分助燃作用。
    \item 煤粉与空气进入炉膛燃烧，底部灰渣经排渣装置排出。
  \end{itemize}
\end{column}
\begin{column}{0.48\textwidth}
  \begin{keybox}{一次风与二次风}
  \small \textbf{一次风}：输送煤粉并参与初期燃烧；\textbf{二次风}：从燃烧器补充燃烧所需空气，组织混合与燃尽。
  \end{keybox}
\end{column}
\end{columns}
\end{frame}

%! src=18 kind=orig title=风烟系统与余热回收
\begin{frame}[t,shrink]{\src{18}{风烟系统与余热回收}}
课件流程：
\begin{center}
\small 烟气 $\rightarrow$ 过热器 $\rightarrow$ 再热器 $\rightarrow$ 省煤器 $\rightarrow$ 空气预热器 $\rightarrow$ 除尘器 $\rightarrow$ 引风机 $\rightarrow$ 脱硫脱硝 $\rightarrow$ 烟囱
\end{center}
\begin{columns}[T,onlytextwidth]
\begin{column}{0.53\textwidth}
  \begin{smalltab}
  \begin{tabularx}{\linewidth}{@{}>{\bfseries}p{42mm}X@{}}
    \toprule
    课件示例（\(600\,\mathrm{MW}\)） & 温度\\
    \midrule
    进口烟温 / 排烟温度 & \(1000/128\,\mathrm{^\circ C}\)\\
    过热汽出口 / 省煤器水温 & \(541/310\,\mathrm{^\circ C}\)\\
    一次风 / 二次风温度 & \(284/307\,\mathrm{^\circ C}\)\\
    \bottomrule
  \end{tabularx}
  \end{smalltab}
  \vspace{0.5mm}
  {\fontsize{6.6}{8}\selectfont\textcolor{black!55}{注：表中数值来自课件示例，随容量、煤种、受热面布置变化。}}
\end{column}
\begin{column}{0.45\textwidth}
  \begin{itemize}
    \item 烟气沿流程逐级放热，分别加热蒸汽、给水和空气。
    \item 空预器把烟气余热送回炉膛，提高助燃空气温度。
    \item 降低排烟温度可减小 \(q_2\)，但受酸露点、低温腐蚀和堵灰限制。
    \item 净化设备的具体布置因机组而异，按课件示意记忆。
  \end{itemize}
\end{column}
\end{columns}
\vspace{0.5mm}
\begin{amberbox}{不要只追求“排烟温度越低越好”}
\small 排烟温度过低会增加空预器面积、引风机电耗和低温腐蚀风险；工程上要综合传热、材料、环保与运行成本。
\end{amberbox}
\end{frame}

%! src=19 kind=orig title=煤粉与风烟系统的耦合
\begin{frame}[t,shrink]{\src{19}{煤粉与风烟系统的耦合}}
\begin{center}
\begin{tikzpicture}[x=1mm,y=1mm,scale=0.86,transform shape]
  \node[box,text width=25mm,minimum height=12mm] (furn) at (65,18) {\textbf{炉膛}\\燃烧与传热};
  \node[box,text width=22mm] (coal) at (15,30) {煤粉};
  \node[box,text width=22mm] (pa) at (15,18) {一次风};
  \node[box,text width=22mm] (sa) at (15,6) {二次风};
  \node[box,text width=25mm] (steam) at (118,30) {蒸汽去汽轮机};
  \node[boxg,text width=25mm] (gas) at (118,18) {烟气去尾部烟道};
  \node[boxg,text width=20mm] (ash) at (118,6) {灰渣};
  \draw[arr] (coal) -- (furn);
  \draw[arr] (pa) -- (furn);
  \draw[arr] (sa) -- (furn);
  \draw[arr] (furn) -- (steam);
  \draw[arr] (furn) -- (gas);
  \draw[arr] (furn) -- (ash);
  \node[lab,text=camber] at (65,4) {\(q_2,q_3,q_4\) 等损失};
\end{tikzpicture}
\end{center}
\begin{columns}[T,onlytextwidth]
\begin{column}{0.49\textwidth}
  \begin{itemize}
    \item \textbf{干燥输粉}：热风把煤粉送入炉膛，并降低煤粉水分。
    \item \textbf{供氧燃烧}：一次风、二次风共同提供氧气并组织混合。
    \item \textbf{传热}：烟气把热量传给水冷壁、过热器、再热器、省煤器和空预器。
  \end{itemize}
\end{column}
\begin{column}{0.48\textwidth}
  \begin{keybox}{风量不是越大越好}
  \small 风量不足会增加 \(q_3\)、\(q_4\)；风量过大又增加烟气量和 \(q_2\)。运行中要同时看燃尽率与排烟损失。
  \end{keybox}
\end{column}
\end{columns}
\end{frame}

%! src=20 kind=orig title=汽水系统的状态变化
\begin{frame}[t,shrink]{\src{20}{汽水系统的状态变化}}
\begin{columns}[T,onlytextwidth]
\begin{column}{0.53\textwidth}
  \textbf{主蒸汽流程}\par\vspace{0.6mm}
  {\small 给水 $\rightarrow$ 省煤器 $\rightarrow$ 饱和水 $\rightarrow$ 水冷壁蒸发 $\rightarrow$ 饱和蒸汽 $\rightarrow$ 过热器 $\rightarrow$ 过热蒸汽 $\rightarrow$ 汽轮机高压缸}
  \par\vspace{2mm}
  \textbf{再热流程}\par\vspace{0.6mm}
  {\small 高压缸排汽 $\rightarrow$ 再热器 $\rightarrow$ 再热蒸汽 $\rightarrow$ 中压缸}
  \par\vspace{2mm}
  \begin{itemize}
    \item 省煤器主要预热，水冷壁主要蒸发，过热器主要过热。
    \item 过热器入口是饱和蒸汽；再热器入口是高压缸排汽，两者来源不同。
    \item 本图按亚临界汽包锅炉理解；超临界直流锅炉没有汽包。
  \end{itemize}
\end{column}
\begin{column}{0.44\textwidth}
\begin{tikzpicture}[x=5.4mm,y=5.4mm]
  \draw[axis] (0,0) -- (10.3,0) node[below,lab] {\(s\)};
  \draw[axis] (0,0) -- (0,7.2) node[left,lab] {\(T\)};
  \draw[thick,cprimary!70] (2.4,0.7) .. controls (3.3,5.5) and (4.0,6.1) .. (4.4,5.9) .. controls (5.4,4.0) and (6.0,2.0) .. (6.5,0.7);
  \draw[thick,cblue] (0.9,1.6) -- (2.7,2.2) -- (5.2,2.2) -- (7.0,5.2) -- (7.9,3.0) -- (9.3,4.8) -- (9.9,1.4);
  \draw[dashline] (2.7,0) -- (2.7,2.2);
  \draw[dashline] (5.2,0) -- (5.2,2.2);
  \foreach \x/\y/\n in {0.9/1.6/1,2.7/2.2/2,5.2/2.2/3,7.0/5.2/4,7.9/3.0/5,9.3/4.8/6,9.9/1.4/7}{
    \fill[cprimary] (\x,\y) circle (0.16);
    \node[font=\fontsize{5.2}{6}\selectfont,text=white] at (\x,\y) {\n};
  }
  \node[lab,text=black!50] at (3.9,1.25) {湿蒸汽区};
\end{tikzpicture}
\figcap{定性 \(T\)-\(s\) 图：1 给水；2 饱和水；3 饱和汽；4 过热汽；5 高压缸排汽；6 再热汽；7 低压排汽。}
\end{column}
\end{columns}
\end{frame}

%! src=20 kind=supp title=过热、再热与给水预热的区别
\begin{frame}[t,shrink]{\srcsupp{20}{过热、再热与给水预热的区别}}
\begin{smalltab}
\begin{tabularx}{\linewidth}{@{}>{\bfseries}p{21mm}p{39mm}X@{}}
  \toprule
  项目 & 入口来源 & 加热目的\\
  \midrule
  给水预热 & 凝结水 / 给水系统 & 回收烟气余热，降低后续加热负担\\
  蒸汽过热 & 汽包分离后的饱和蒸汽 & 提高主蒸汽温度与做功能力\\
  蒸汽再热 & 高压缸已经膨胀做功的排汽 & 再次升温后进入中压缸继续膨胀\\
  \bottomrule
\end{tabularx}
\end{smalltab}
\begin{columns}[T,onlytextwidth]
\begin{column}{0.53\textwidth}
  \begin{itemize}
    \item 再热的前提是蒸汽\textbf{已经在汽轮机内膨胀过一次}。
    \item 再热蒸汽的压力低于主蒸汽，不能把两股流量和焓差直接合并。
    \item 给水预热发生在锅炉入口之前，不属于“蒸汽过热”。
  \end{itemize}
\end{column}
\begin{column}{0.44\textwidth}
  \begin{amberbox}{边界}
  \small 计算锅炉有效吸热时，要分别列出主蒸汽和再热蒸汽的焓升；再热器漏泄、减温水等支流会使实际流量与主蒸汽流量不同。
  \end{amberbox}
\end{column}
\end{columns}
\end{frame}

%! src=21 kind=orig title=汽包锅炉的汽水循环
\begin{frame}[t,shrink]{\src{21}{汽包锅炉的汽水循环}}
\begin{columns}[T,onlytextwidth]
\begin{column}{0.49\textwidth}
  \begin{itemize}
    \item \textbf{给水进入}：凝结水与补充水 $\rightarrow$ 给水泵 $\rightarrow$ 省煤器 $\rightarrow$ 汽包。
    \item \textbf{内部循环}：汽包 $\rightarrow$ 下降管 $\rightarrow$ 下联箱 $\rightarrow$ 水冷壁 $\rightarrow$ 上联箱 $\rightarrow$ 汽包。
    \item \textbf{蒸汽送出}：汽包分离出的饱和蒸汽 $\rightarrow$ 过热器 $\rightarrow$ 汽轮机。
    \item 下降管中主要是水；水冷壁出口为汽水混合物。
  \end{itemize}
\end{column}
\begin{column}{0.48\textwidth}
\begin{tikzpicture}[x=1mm,y=1mm,scale=0.9,transform shape]
  \draw[cprimary,thick] (6,5) rectangle (54,42);
  \draw[cprimary,thick] (30,42) circle (4.5);
  \draw[cprimary,thick] (13,42) -- (13,11) -- (47,11) -- (47,42);
  \draw[arr] (13,30) -- (13,13);
  \draw[arr] (47,13) -- (47,40);
  \node[lab,anchor=south] at (30,47) {汽包};
  \node[lab,anchor=east] at (10,26) {下降管};
  \node[lab,anchor=west] at (50,26) {水冷壁};
  \node[lab] at (30,7) {下联箱};
\end{tikzpicture}
\end{column}
\end{columns}
\vspace{0.8mm}
\begin{amberbox}{读图重点}
\small 分清三条关系：外部给水、内部循环水、送出的蒸汽。不要把下降管、上升管和过热器串成一条“单程路线”。
\end{amberbox}
\end{frame}

%! src=22 kind=orig title=煤的分析方法与成分基准
\begin{frame}[t,shrink]{\src{22}{煤的分析方法与成分基准}}
\begin{columns}[T,onlytextwidth]
\begin{column}{0.50\textwidth}
  \begin{itemize}
    \item \textbf{元素分析}：碳、氢、氧、氮、硫等元素，以及灰分、水分。
    \item \textbf{工业分析}：水分、灰分、挥发分、固定碳。
    \item 两种分析的划分依据不同；固定碳不等于全部元素碳。
  \end{itemize}
  \vspace{0.8mm}
  \begin{smalltab}
  \begin{tabularx}{\linewidth}{@{}>{\bfseries}p{18mm}p{10mm}X@{}}
    \toprule
    基准 & 下标 & 含义\\
    \midrule
    收到基 & \(\mathrm{ar}\) & 入炉煤的全部成分\\
    空气干燥基 & \(\mathrm{ad}\) & 风干状态，去除部分外部水\\
    干燥基 & \(\mathrm{d}\) & 去掉全部水分\\
    干燥无灰基 & \(\mathrm{daf}\) & 去掉全部水分和灰分\\
    \bottomrule
  \end{tabularx}
  \end{smalltab}
\end{column}
\begin{column}{0.47\textwidth}
\begin{tikzpicture}[x=1mm,y=1mm,scale=0.9,transform shape]
  \draw[fill=cprimary!8,draw=cprimary!50,rounded corners=2pt] (0,0) rectangle (55,40);
  \draw[fill=cprimary!14,draw=cprimary!50,rounded corners=2pt] (4,4) rectangle (51,36);
  \draw[fill=cprimary!22,draw=cprimary!50,rounded corners=2pt] (8,8) rectangle (47,32);
  \draw[fill=cprimary!32,draw=cprimary!50,rounded corners=2pt] (12,12) rectangle (43,28);
  \node[font=\fontsize{6.0}{7.2}\selectfont,anchor=north west] at (1,39) {收到基 ar};
  \node[font=\fontsize{6.0}{7.2}\selectfont,anchor=north west] at (5,35) {空气干燥基 ad};
  \node[font=\fontsize{6.0}{7.2}\selectfont,anchor=north west] at (9,31) {干燥基 d};
  \node[font=\fontsize{6.0}{7.2}\selectfont,anchor=north west] at (13,27) {干燥无灰基 daf};
  \node[font=\fontsize{5.8}{7}\selectfont,text=black!55] at (27,18) {水分、灰分逐层扣除};
\end{tikzpicture}
\figcap{基准越“干”、越“无灰”，分母越小，同一组分的百分数通常越高。}
\end{column}
\end{columns}
\end{frame}

%! src=22 kind=supp title=煤质基准转换的分母
\begin{frame}[t,shrink]{\srcsupp{22}{煤质基准转换：分母到底变成什么}}
对同一煤样中某个非水、非灰的质量组分 \(X\)，以质量百分数表示：
$$
X_{\mathrm{d}}
=X_{\mathrm{ar}}\frac{100}{100-M_{\mathrm{ar}}},
\qquad
X_{\mathrm{daf}}
=X_{\mathrm{ar}}\frac{100}{100-M_{\mathrm{ar}}-A_{\mathrm{ar}}}
$$
\begin{columns}[T,onlytextwidth]
\begin{column}{0.50\textwidth}
  \begin{itemize}
    \item \(M_{\mathrm{ar}}\)、\(A_{\mathrm{ar}}\)：收到基水分、灰分的质量百分数。
    \item 去掉水或灰后，所考察组分的质量不变，但计算百分数的总质量减小。
    \item 因此同一组分换算到干燥基或干燥无灰基时，百分数通常升高。
  \end{itemize}
\end{column}
\begin{column}{0.47\textwidth}
  \begin{amberbox}{适用边界}
  \small 这是\textbf{质量组成}的换算。发热量还涉及水分汽化、定义基准和氧弹条件，不能把同一公式无条件套到发热量上。
  \end{amberbox}
  \vspace{0.7mm}
  \begin{keybox}{检查方法}
  \small 先问“分母里还剩哪些成分”，再检查换算后各组分之和是否为 \(100\%\)。
  \end{keybox}
\end{column}
\end{columns}
\end{frame}

%! src=23 kind=orig title=锅炉热平衡的六项输出
\begin{frame}[t,shrink]{\src{23}{锅炉热平衡的六项输出}}
$$
Q_{\mathrm{f}}=Q_1+Q_2+Q_3+Q_4+Q_5+Q_6
$$
\begin{columns}[T,onlytextwidth]
\begin{column}{0.52\textwidth}
  \begin{center}
  \begin{tikzpicture}[x=1mm,y=1mm,scale=0.9,transform shape]
    \node[box,text width=22mm,minimum height=14mm] (qf) at (18,18) {输入\\\(Q_{\mathrm{f}}\)};
    \node[boxb,text width=25mm] (q1) at (70,34) {\(Q_1\) 有效利用};
    \node[boxg,text width=25mm] (q2) at (70,26) {\(Q_2\) 排烟};
    \node[boxg,text width=25mm] (q3) at (70,18) {\(Q_3\) 未燃气体};
    \node[boxg,text width=25mm] (q4) at (70,10) {\(Q_4\) 未燃固体};
    \node[boxg,text width=25mm] (q5) at (70,2) {\(Q_5\) 散热};
    \node[boxg,text width=25mm] (q6) at (70,-6) {\(Q_6\) 灰渣物理热};
    \draw[arr] (qf.east) -- (q1.west);
    \draw[arr] (qf.east) -- (q2.west);
    \draw[arr] (qf.east) -- (q3.west);
    \draw[arr] (qf.east) -- (q4.west);
    \draw[arr] (qf.east) -- (q5.west);
    \draw[arr] (qf.east) -- (q6.west);
  \end{tikzpicture}
  \end{center}
\end{column}
\begin{column}{0.45\textwidth}
  \begin{smalltab}
  \begin{tabularx}{\linewidth}{@{}>{\bfseries}p{11mm}X@{}}
    \toprule
    项目 & 含义\\
    \midrule
    \(Q_1\) & 向工质传递的有效利用热\\
    \(Q_2\) & 排烟及其携带飞灰带走的热\\
    \(Q_3\) & 未燃尽可燃气体带走的热\\
    \(Q_4\) & 飞灰、灰渣中未燃尽可燃物的热\\
    \(Q_5\) & 炉墙向环境散热\\
    \(Q_6\) & 灰渣物理热等其他损失\\
    \bottomrule
  \end{tabularx}
  \end{smalltab}
\end{column}
\end{columns}
\vspace{0.5mm}
\begin{amberbox}{符号约定}
\small \(Q_{\mathrm{f}}\) 是输入热量；本页各 \(Q\) 可以按每千克燃料计，也可以按同一时间段计。若写成热流率，则全部改为 \(\dot Q\)。
\end{amberbox}
\end{frame}

%! src=23 kind=supp title=把热平衡统一到同一基准
\begin{frame}[t,shrink]{\srcsupp{23}{把热平衡统一到同一时间或同一质量基准}}
输入、有效利用热量和损失可以按同一时间段计，也可全部改为热流率：
$$
\dot Q_{\mathrm{f}}=\dot Q_1+\sum_{j=2}^{6}\dot Q_j,
\qquad
\eta_{\mathrm{b}}=\frac{\dot Q_1}{\dot Q_{\mathrm{f}}}
$$
\begin{columns}[T,onlytextwidth]
\begin{column}{0.51\textwidth}
  \begin{itemize}
    \item 若简化输入为燃料低位发热量：\(\dot Q_{\mathrm{f}}\approx\dot m_{\mathrm{f}}H_{\mathrm{u}}\)。
    \item 空气显热、燃料显热、雾化蒸汽带入热量是否计入，应服从题目边界。
    \item \(Q\) 用 \(\mathrm{kJ}\)，\(\dot Q\) 用 \(\mathrm{kW}\)，不能把热量与功率直接相除。
  \end{itemize}
\end{column}
\begin{column}{0.46\textwidth}
  \begin{keybox}{百分数的两种写法}
  \small 若 \(q_j=Q_j/Q_{\mathrm{f}}\) 为小数，则 \(\eta_{\mathrm{b}}=q_1\)；若 \(q_j\) 已经写成百分数，则不能再乘 \(100\%\)。原课件把两者写在同一行，读式子时要先看 \(q_1\) 的定义。
  \end{keybox}
\end{column}
\end{columns}
\end{frame}

%! src=24 kind=orig title=正平衡、反平衡与热损失
\begin{frame}[t,shrink]{\src{24}{正平衡、反平衡与热损失}}
统一采用小数形式 \(q_j=Q_j/Q_{\mathrm{f}}\)：
$$
\eta_{\mathrm{b}}=q_1
=1-\left(q_2+q_3+q_4+q_5+q_6\right)
$$
\begin{columns}[T,onlytextwidth]
\begin{column}{0.48\textwidth}
  \begin{itemize}
    \item \textbf{正平衡}：测有效利用热量，求 \(\eta_{\mathrm{b}}=Q_1/Q_{\mathrm{f}}\)。
    \item \textbf{反平衡}：测各项损失，再用 1 减去损失之和。
    \item 两种方法应相互校验；差值过大时先查测量与边界。
  \end{itemize}
\end{column}
\begin{column}{0.48\textwidth}
  \begin{keybox}{课件经验值（有条件的示例）}
  \small
  排烟温度升高 \(10\sim15\,\mathrm{^\circ C}\)，\(q_2\) 约增加 1 个百分点；\\
  煤粉锅炉 \(q_3\approx0\)，\(q_4\approx0.5\%\sim1\%\)；\\
  \(D>900\,\mathrm{t/h}\) 时，课件列 \(q_5\approx0.2\%\)。
  \end{keybox}
\end{column}
\end{columns}
\end{frame}

%! src=24 kind=supp title=从损失项反推运行检查
\begin{frame}[t,shrink]{\srcsupp{24}{从损失项反推应检查的运行因素}}
\begin{smalltab}
\begin{tabularx}{\linewidth}{@{}>{\bfseries}p{10mm}p{35mm}X@{}}
  \toprule
  损失 & 主要检查对象 & 可能的因果链\\
  \midrule
  \(q_2\) & 排烟温度、烟气量 & 受热面积灰、漏风、过量空气系数过大\\
  \(q_3\) & 未燃尽可燃气体 & 燃烧组织不良、局部缺氧、温度过低\\
  \(q_4\) & 飞灰、炉渣含碳量 & 煤粉过粗、燃尽时间不足、配风不当\\
  \(q_5\) & 炉墙保温、机组负荷 & 保温损坏；低负荷时固定散热占比升高\\
  \(q_6\) & 灰渣温度、灰分与排渣方式 & 灰渣显热未回收；干排渣与湿排渣不同\\
  \bottomrule
\end{tabularx}
\end{smalltab}
\begin{columns}[T,onlytextwidth]
\begin{column}{0.51\textwidth}
  \begin{itemize}
    \item 先找损失对应的\textbf{能量出口}，再判断哪个运行参数可能起作用。
    \item 各项损失相互耦合：加大风量可能降低 \(q_3,q_4\)，却提高 \(q_2\)。
    \item 课件末条“\(q_6\approx0\)”的判据在投影页上未给全单位，使用前应查教材原式。
  \end{itemize}
\end{column}
\begin{column}{0.46\textwidth}
  \begin{amberbox}{不要背成单向口诀}
  \small “降低排烟温度”并不总是净收益；还要考虑酸露点、空预器面积、引风机电耗和低温腐蚀。
  \end{amberbox}
\end{column}
\end{columns}
\end{frame}
